\chapter{Fecal Fat Test}

\section*{What Is This Test?}
This stool test measures how much fat remains in the stool instead of being absorbed. A high amount can be a sign that the pancreas, small intestine, or bile system is not helping the body absorb food properly. This may cause oily or floating stools, diarrhea, weight loss, or vitamin deficiencies.

\section*{Alternative Names}
\begin{itemize}
  \item Stool Fat Test
  \item Faecal Fat
  \item Quantitative Fecal Fat
  \item Sudan Stain for Stool Fat
  \item Steatorrhea Test
\end{itemize}
\section*{How It Is Performed}
A quick test may look for fat droplets in one stool sample. A more exact test may require collecting all stool for a set time, often 72 hours, while eating the amount of fat specified by the laboratory. The laboratory then measures the total fat in the collected stool.

\section*{Why It Is Ordered}
\begin{itemize}
  \item Investigation of chronic diarrhea, bulky or oily stools, weight loss, and suspected steatorrhea
  \item Evaluation of suspected pancreatic exocrine insufficiency or small-bowel malabsorption
  \item Assessment of persistent nutritional deficiencies when routine evaluation is inconclusive
\end{itemize}

\section*{Preparation, Risks, and Limitations}
There is no needle or procedure risk, but collecting stool for several days can be inconvenient. Follow the diet and collection instructions carefully. Missing part of the stool, collecting urine or water with it, or eating less fat than instructed can make the result unreliable. A high result does not show the cause; other tests may be needed.
\section*{Specimen, Preparation, and Safety}
Use the laboratory's container and keep urine, water, and cleaning products out of the sample. Follow the requested diet and collection period. If any stool is missed, tell the laboratory because the test may need to be repeated.

\section*{Understanding Results and Follow-up}
The report may say whether fat was seen or give the amount of fat collected per day. A high result supports poor fat absorption but does not identify the cause. Follow-up may include a pancreatic stool test, celiac blood tests, liver tests, or a nutrition assessment. The clinician also considers your symptoms, diet, and whether the collection was complete.

\section*{Regulatory Status and Authorization Date}
This chapter describes a test category, not a specific commercial product. FDA, CDSCO, EU IVDR, and MHRA approval or registration dates therefore cannot be assigned without the assay manufacturer, product name, instrument, and intended use. For laboratory-developed tests, an individual agency approval date may not exist. See the Preface for the interpretation of regulatory status.

\clearpage
